\documentclass[11pt]{article}
\usepackage[T1]{fontenc}
\usepackage{amsmath,amssymb}
\begin{document}
\section*{EPPA--III, Lemma 3.18: a one-automorphism Petersen primitivity certificate}
The author manuscript is frozen. This is a proposed local proof
of the remaining primitive/imprimitive distinction in the matching case.

Use the canonical matching model on
\(V=\{0,1\}\times\mathbb Z/5\mathbb Z\).
The bottom layer is \(C_5\), the top is
\(\overline{C_5}\), and the cross edges form a perfect
matching. Number the vertices by \(0,\ldots,4\)
on the bottom and \(5,\ldots,9\) on the top.

A direct ten-by-ten adjacency check shows that
\[
   \sigma=(1\ 4\ 5)(2\ 3\ 8\ 6\ 9\ 7)
\]
(with \(0\) fixed) is an automorphism of this
matching graph. Its orbits are of sizes \(1,3,6\).
Any imprimitivity block containing \(0\) is fixed
setwise by \(\sigma\), because \(\sigma(0)=0\) and
images of blocks are equal or disjoint.
Thus such a block must be a union of the three
orbits and its size is one of
\[
                   1,\ 4,\ 7,\ 10.
\]
In any transitive action on ten vertices, a
nontrivial invariant partition has block size
\(2\) or \(5\). Neither occurs, so the Petersen
matching graph is primitive.

This is considerably shorter than enumerating
the \(120\) graph automorphisms or proving the full
permutation-group identification with \(S_5\).
It requires only one explicit automorphism,
one calculation on at most \(2^{10}=1024\) subsets,
and the ordinary equal-fibre cardinality lemma
already formalized for Section 3.

\paragraph{Lean target.}
The proposed module
\texttt{CycleFivePetersenPrimitivity.lean}
uses an explicit computable permutation of
the ten canonical vertices, proves its graph
automorphism property, and certifies by finite
kernel reduction that no root-containing invariant
subset can have size two or five.
It then composes this certificate with the
already verified \texttt{transitive\_block\_fibres\_equal}
and \texttt{uniform\_partition\_card\_formula}.
The resulting theorem
\texttt{petersen\_no\_nontrivial\_invariant\_partition}
assumes vertex-transitivity and excludes a nontrivial
invariant block partition.

\paragraph{Still needed after this certificate.}
Transport any given invariant block map of the actual
host along the canonical graph equivalence from
\texttt{cycleFive\_eppa\_singleton\_host\_two\_matching\_graphs}.
This leaves the host exclusion as a short corollary and
then permits the full imprimitive matching case of
Lemma~3.18 to be assembled.

\paragraph{Suggested source-local insertion.}
The Petersen alternative is primitive: in the
matching coordinates the automorphism
\((1\ 4\ 5)(2\ 3\ 8\ 6\ 9\ 7)\) fixes vertex \(0\).
A block containing \(0\) has order \(1,4,7\), or
\(10\), whereas transitivity on ten vertices
forces any nontrivial block to have order \(2\) or \(5\).
\end{document}
